\documentclass[12pt]{article}
\usepackage[T1]{fontenc}
\usepackage[utf8]{inputenc}
\usepackage{lmodern}
\usepackage{textcomp}
\usepackage{enumitem}
\usepackage{geometry}
\geometry{margin=1in}
\begin{document}
\begin{center}
Answer Key - Chemistry - Undergraduate Level Assessment 15 \
\small (Answers are ordered exactly as in the question paper.)
\end{center}

\section*{Multiple Choice}
\begin{enumerate}[label=\arabic*.]
\item \textbf{Answer:} A
\\ \textit{Options:} A) B$_{12}$ icosahedra; B) B$_{8}$ cubes; C) B$_{6}$ octahedra; D) B$_{4}$ tetrahedra
\\ \textit{Note:} The correct answer is B$_{12}$ icosahedra because the principal allotropes of elemental boron, such as alpha-boron and beta-boron, are characterized by their unique three-dimensional networks formed by interconnected B$_{12}$ icosahedral units, which are essential to their solid-state structures.
\item \textbf{Answer:} C
\\ \textit{Options:} A) 3; B) 2; C) 1; D) 0
\\ \textit{Note:} The orbital angular momentum quantum number, l, of 1 corresponds to the p-orbital, which is the highest energy level electron in ground-state aluminum (atomic number 13), making it the most easily removed during ionization.
\item \textbf{Answer:} C
\\ \textit{Options:} A) all molecular bonds absorb IR radiation; B) IR peak intensities are related to molecular mass; C) most organic functional groups absorb in a characteristic region of the IR spectrum; D) each element absorbs at a characteristic wavelength
\\ \textit{Note:} The correct answer is based on the fact that different organic functional groups have unique vibrational frequencies that correspond to specific regions of the infrared spectrum, allowing for the identification of these groups in a molecule's structure.
\item \textbf{Answer:} D
\\ \textit{Options:} A) racemic mixture; B) single pure enantiomer; C) mixture of two diastereomers in equal amounts; D) meso compound
\\ \textit{Note:} The reduction of D-xylose with NaBH 4 results in a meso compound because the product exhibits internal symmetry and has multiple stereocenters, leading to a compound that is optically inactive despite having chiral centers.
\item \textbf{Answer:} D
\\ \textit{Options:} A) Ammonium hydroxide; B) Sulfuric acid; C) Acetic acid; D) Potassium hydrogen phthalate
\\ \textit{Note:} Potassium hydrogen phthalate is a primary standard because it is a stable, non-hygroscopic substance that can be accurately weighed and dissolved to prepare standard solutions for the precise standardization of bases in titrations.
\end{enumerate}

\section*{True / False}
\begin{enumerate}[label=\arabic*.]
\setcounter{enumi}{5}
\item \textbf{Answer:} True
\\ \textit{Options:} A) True; B) False
\\ \textit{Note:} At standard temperature and pressure, bromine (Br$_{2}$) exists as a red-brown liquid due to its molecular structure and intermolecular forces, which allow it to be volatile at room temperature.
\item \textbf{Answer:} False
\\ \textit{Options:} A) True; B) False
\\ \textit{Note:} The statement is false because BeCl 2 acts as a Lewis acid primarily due to the presence of an incomplete octet around the beryllium atom, making it more electron-deficient, while the higher electronegativity of chlorine in SCl 2 does not influence its ability to act as a Lewis acid in the same manner.
\item \textbf{Answer:} True
\\ \textit{Options:} A) True; B) False
\\ \textit{Note:} The statement is true because protons and neutrons, as nucleons within the atomic nucleus, do not exhibit orbital angular momentum like electrons, but rather possess intrinsic spin angular momentum.
\item \textbf{Answer:} False
\\ \textit{Options:} A) True; B) False
\\ \textit{Note:} Tetramethylsilane is used as a reference in NMR spectroscopy because its protons are highly shielded, leading to a distinct and consistent peak that serves as a standard for chemical shift measurements.
\item \textbf{Answer:} False
\\ \textit{Options:} A) True; B) False
\\ \textit{Note:} The statement is false because, in addition to symmetric stretching, the normal modes of a carbon dioxide molecule that are infrared-active also include bending modes, specifically the asymmetric bending mode.
\end{enumerate}

\section*{Long Form Response}
\begin{enumerate}[label=\arabic*.]
\setcounter{enumi}{10}
\item \textbf{Answer:} The limiting high-temperature molar heat capacity at constant volume (C\_V) of a gas-phase diatomic molecule is primarily influenced by the degrees of freedom available to the molecule, which include translational, rotational, and vibrational motions. At high temperatures, diatomic molecules exhibit three translational and two rotational degrees of freedom, contributing a total of 5 R to C\_V. However, as vibrational modes begin to be significantly excited at elevated temperatures, each vibrational mode adds an additional R/2 to C\_V, leading to a theoretical maximum of 7 R for a diatomic molecule. Yet, the vibrational contribution becomes significant only at very high temperatures, so at typical high temperatures, C\_V approaches 5 R. Thus, the limiting value of C\_V is often expressed as 3.5 R due to the predominance of translational and rotational contributions before vibrational modes are fully excited. This relationship highlights the interplay between molecular degrees of freedom and thermodynamic properties in diatomic gases.
\\ \textit{Note:} correct because it accurately identifies the contributions of translational, rotational, and vibrational degrees of freedom to the molar heat capacity of diatomic molecules at high temperatures, as well as the significance of vibrational excitation in determining C\_V.
\item \textbf{Answer:} The Heisenberg uncertainty principle has profound implications for a quantum mechanical particle confined in the lowest energy level of a one-dimensional box. According to this principle, it is impossible to precisely determine both the position and momentum of the particle simultaneously; the more accurately we know one, the less accurately we can know the other. In the lowest energy state, the particle is described by a wavefunction that is spread out over the entire box, indicating a high degree of uncertainty in its position. Consequently, this spatial uncertainty leads to a correspondingly high uncertainty in momentum, reflecting the wave-particle duality fundamental to quantum mechanics. Thus, the Heisenberg uncertainty principle emphasizes the inherent limitations in measuring quantum properties, illustrating a departure from classical mechanics where both position and momentum can be known exactly.
\\ \textit{Note:} correct because it accurately describes how the Heisenberg uncertainty principle limits our ability to precisely measure both the position and momentum of a quantum particle, which is illustrated by the wavefunction of the particle being spread out in the lowest energy state of a one-dimensional box.
\end{enumerate}

\vfill
\noindent\textit{End of Answer Key}
\end{document}